%% ============================================================
%%  thesis_info.tex — ЄДИНИЙ ФАЙЛ ДЛЯ ЗАПОВНЕННЯ
%% ============================================================

%% ══════════════════════════════════════════════════════
%%  БЛОК 1 — Титульна сторінка
%% ══════════════════════════════════════════════════════

%% Повна назва роботи
\renewcommand{\ThesisTitle}{Вкажіть назву роботи}

%% Прізвище Ім'я По-батькові
\renewcommand{\ThesisAuthor}{Вкажіть прізвище ім'я по-батькові}

%% Посада, науковий ступінь, Прізвище І.П.
\renewcommand{\ThesisSupervisor}{Вкажіть наукового керівника}

%% Офіційна назва кафедри
\renewcommand{\ThesisDept}{Вкажіть кафедру}

%% Офіційна назва факультету або інституту
\renewcommand{\ThesisFaculty}{Вкажіть факультет}

%% Офіційна назва університету
\renewcommand{\ThesisUniversity}{Вкажіть університет}

%% Код та назва спеціальності: 122 "Комп'ютерні науки"
\renewcommand{\ThesisSpecialty}{Вкажіть спеціальність}

%% Місто та рік захисту
\renewcommand{\ThesisCity}{Київ}
\renewcommand{\ThesisYear}{2025}

%% ══════════════════════════════════════════════════════
%%  БЛОК 2 — Реферат та вступ (генеруються автоматично)
%% ══════════════════════════════════════════════════════

%% Ключові слова ВЕЛИКИМИ ЛІТЕРАМИ через кому (5–8 слів)
\renewcommand{\ThesisKeywords}{ВКАЖІТЬ КЛЮЧОВІ СЛОВА}

%% Мета роботи — починати з інфінітива ("розробити...", "дослідити...")
\renewcommand{\ThesisMeta}{вкажіть мету роботи}

%% Об'єкт дослідження — процес, явище або система
\renewcommand{\ThesisObject}{вкажіть об'єкт дослідження}

%% Предмет дослідження — методи, моделі, характеристики
\renewcommand{\ThesisSubject}{вкажіть предмет дослідження}

%% Методи дослідження — перерахування через кому
\renewcommand{\ThesisMethods}{вкажіть методи дослідження}

%% Наукова новизна — що вперше запропоновано або удосконалено
\renewcommand{\ThesisNovelty}{вкажіть наукову новизну}

%% Практичне значення отриманих результатів
\renewcommand{\ThesisPractical}{вкажіть практичне значення}

%% ══════════════════════════════════════════════════════
%%  БЛОК 3 — Необов'язкові поля
%%  Якщо не потрібно — залиште порожніми {}
%% ══════════════════════════════════════════════════════

%% Де доповідались результати (конференції, семінари)
\renewcommand{\ThesisAprobation}{}

%% Публікації за темою роботи
\renewcommand{\ThesisPublications}{}

%% ══════════════════════════════════════════════════════
%%  СТАТИСТИКА (рисунки/таблиці/джерела/додатки)
%%
%%  Рахується АВТОМАТИЧНО — нічого вписувати не треба.
%%  Значення з'являються після 2-ї компіляції (latexmk робить це сам).
%%  Якщо у рефераті стоять "?" — просто перекомпілюйте ще раз.
%% ══════════════════════════════════════════════════════